\section{Contributi e obiettivi della tesi}
\label{sec:contributi_obiettivi}

L'obiettivo principale di questo lavoro non è la mera riduzione
del consumo energetico medio di una piattaforma serverless, bensì
la progettazione e realizzazione di un'estensione energy-aware della
piattaforma Serverledge~\cite{Russo2023_Serverledge} in grado di
operare sotto vincoli energetici espliciti e di adattare le decisioni
di esecuzione in funzione di informazioni energetiche aggiornate
a runtime.

In particolare, il lavoro si propone di:

\begin{itemize}

  \item introdurre il concetto di \emph{funzione logica} associata
  a un insieme di varianti implementative alternative, ciascuna
  descritta da un profilo energetico che ne quantifica il costo
  atteso per invocazione, il costo di cold start e, ove applicabile,
  la stima dell'errore introdotto dall'approssimazione;

  \item integrare un sistema di monitoraggio energetico a livello
  di container basato su Kepler, strumento che stima il consumo
  energetico a granularità di processo nei nodi Kubernetes e
  lo esporta tramite Prometheus, consentendo la raccolta di
  misure energetiche reali durante l'esecuzione delle funzioni;

  \item progettare e implementare un meccanismo di aggiornamento
  dinamico dei profili energetici delle varianti, realizzato tramite
  un worker periodico che legge le misure cumulative da Kepler,
  calcola il consumo per singola invocazione e aggiorna i profili
  archiviati in etcd, realizzando un ciclo di feedback adattivo;

  \item progettare e implementare uno scheduler energeticamente
  consapevole, ispirato concettualmente ai meccanismi di controllo
  basati su budget energetico proposti in letteratura, che al momento
  dell'invocazione seleziona la variante più appropriata in funzione
  di:
    \begin{itemize}
      \item il budget energetico massimo espresso dall'utente
            (\texttt{MaxEnergyJoule});
      \item il consenso esplicito dell'utente all'utilizzo di
            varianti approssimate (\texttt{AllowApprox});
      \item le stime energetiche aggiornate disponibili a runtime;
    \end{itemize}
  rispettando il vincolo che, in assenza di consenso esplicito,
  vengano selezionate esclusivamente varianti esatte;

  \item valutare sperimentalmente il compromesso tra consumo
  energetico e accuratezza del risultato su un insieme di funzioni
  rappresentative, dimostrando la fattibilità e l'efficacia
  di un modello di esecuzione adattivo per il paradigma
  Function-as-a-Service nel continuum Edge--Cloud.

\end{itemize}

Il contributo complessivo della tesi consiste nell'estensione
di una piattaforma FaaS open-source con un meccanismo di esecuzione
adattivo rispetto a vincoli energetici dinamici, introducendo
un compromesso controllato tra energia consumata e accuratezza
del risultato, condizionato al consenso esplicito dell'utente
che effettua l'invocazione.
